% GENERATED by scripts/make_waste_tier_table.py - do not hand-edit;
% a test asserts this file matches a fresh regeneration, and the generator
% asserts every parameter value it quotes against data/waste_params.csv.
% ===========================================================================
\section{Supplementary Table: tier status by component}
\label{app:tiers}
% ===========================================================================

Tier status differs by component. The three categories share one activity
basis -- the district population panel -- but carry different levels of
country-specific characterisation, and two pieces of the sector are not on
the IPCC tier ladder at all. One word for the whole sector would either
overstate the factors or understate the activity.

\begin{table*}[htbp]
\centering
\scriptsize
\begin{tabular}{@{}llp{0.55\linewidth}@{}}
\toprule
\textbf{Component} & \textbf{Status} & \textbf{Basis} \\
\midrule
4A First Order Decay & Tier 1, country-specific activity & IPCC Vol 5 Ch 3 FOD with default $k$ (0.065), DOC$_f$, $F$ and the uncategorised-site MCF (0.6) over a district census population panel and Pakistani generation rates (0.453 / 0.283 kg/cap/day, Pak-EPA 1996). The activity is country-specific; every decay parameter is a default, and the MCF stands in for a site-type census Pakistan has never published. \\[2pt]
4C2 Open burning & Tier 1, country-specific activity & IPCC Eq 5.7 over the uncollected stream (0.30 urban / 0.70 rural, What a Waste) $\times$ the Box 5.1 burned fraction (0.60), with the Pak-EPA composition supplying carbon and dry matter. The burned fraction has no Pakistani measurement behind it. \\[2pt]
4C2 air-pollutant layer & Field factors (not an IPCC tier) & Andreae 2019 garbage-burning factors on the identical burned mass. The IPCC tier ladder classifies greenhouse-gas estimation; a field-measurement factor set for co-emitted pollutants sits beside it, not on it. \\[2pt]
4D1 CH$_4$ & Tier 1, stratified activity & IPCC Vol 5 Ch 6 defaults (BOD 40 g/cap/day regional, $B_0$ 0.6) over the same population panel stratified by the JMP sanitation split - the stratification is country-specific, the factors are defaults, and the sewer MCF (0.5 central, 0.1 published alongside) is the sector's deciding judgement. \\[2pt]
4D1 N$_2$O & Tier 1 & FAOSTAT protein through the 2019 Refinement's nitrogen chain; the effluent factor's published range spans a factor of 150, the softest number in the sector. \\[2pt]
Provincial allocation & Country-specific allocation (not a tier) & Three distinct district-population-derived proxies (burning, disposal, wastewater load), reconciling exactly to each category's national figure. IPCC tiers classify estimation, not allocation. \\[2pt]
\bottomrule
\end{tabular}
\caption{Tier status by component. Generated; the parameter values quoted
are asserted against \texttt{data/waste\_params.csv} at generation time.}
\label{tab:tierstatus}
\end{table*}
